% =============================================================================
%  O.S.C.A.R. — Documentation technique et guide utilisateur
%  Outil de Suivi des Cours et d'Analyse du Réseau
%  Institut français Italia
% =============================================================================
\documentclass[12pt,a4paper,twoside]{article}

% ── Encodage et langue ──
\usepackage[utf8]{inputenc}
\usepackage[T1]{fontenc}
\usepackage[french]{babel}

% ── Mise en page ──
\usepackage[margin=2.5cm, headheight=15pt]{geometry}
\usepackage{fancyhdr}
\usepackage{parskip}
\usepackage{titlesec}
\usepackage{enumitem}

% ── Couleurs et graphiques ──
\usepackage[dvipsnames,table]{xcolor}
\usepackage{graphicx}
\usepackage{tikz}
\usepackage{tcolorbox}

% ── Tableaux ──
\usepackage{booktabs}
\usepackage{tabularx}
\usepackage{longtable}

% ── Code ──
\usepackage{listings}
\usepackage{verbatim}

% ── Liens ──
\usepackage[colorlinks=true, linkcolor=MidnightBlue, urlcolor=RoyalBlue, citecolor=ForestGreen]{hyperref}

% ── Divers ──
\usepackage{microtype}
\usepackage{csquotes}

% ── Couleurs IFI ──
\definecolor{ifiBlue}{HTML}{1E3A5F}
\definecolor{ifiRed}{HTML}{E63946}
\definecolor{ifiGray}{HTML}{6B7280}
\definecolor{ifiLight}{HTML}{F3F4F6}
\definecolor{codeBg}{HTML}{F8F9FA}

% ── Styles titres ──
\titleformat{\section}
  {\Large\bfseries\color{ifiBlue}}{\thesection}{1em}{}[\titlerule]
\titleformat{\subsection}
  {\large\bfseries\color{ifiBlue!80}}{\thesubsection}{1em}{}
\titleformat{\subsubsection}
  {\normalsize\bfseries\color{ifiBlue!60}}{\thesubsubsection}{1em}{}

% ── En-tête / pied de page ──
\pagestyle{fancy}
\fancyhf{}
\fancyhead[L]{\small\color{ifiGray}O.S.C.A.R. --- Documentation}
\fancyhead[R]{\small\color{ifiGray}Institut fran\c{c}ais Italia}
\fancyfoot[C]{\small\color{ifiGray}\thepage}
\renewcommand{\headrulewidth}{0.4pt}
\renewcommand{\footrulewidth}{0pt}

% ── Style code ──
\lstset{
  backgroundcolor=\color{codeBg},
  basicstyle=\ttfamily\small,
  breaklines=true,
  frame=single,
  rulecolor=\color{ifiGray!30},
  numbers=left,
  numberstyle=\tiny\color{ifiGray},
  keywordstyle=\color{ifiBlue}\bfseries,
  commentstyle=\color{ForestGreen}\itshape,
  stringstyle=\color{ifiRed},
  showstringspaces=false,
  tabsize=4,
}

% ── Boîtes info/attention ──
\newtcolorbox{infobox}[1][]{
  colback=blue!5, colframe=ifiBlue, fonttitle=\bfseries,
  title={#1}, arc=2mm, boxrule=0.5mm
}
\newtcolorbox{warnbox}[1][]{
  colback=red!5, colframe=ifiRed, fonttitle=\bfseries,
  title={#1}, arc=2mm, boxrule=0.5mm
}
\newtcolorbox{tipbox}[1][]{
  colback=green!5, colframe=ForestGreen, fonttitle=\bfseries,
  title={#1}, arc=2mm, boxrule=0.5mm
}

% =============================================================================
\begin{document}

% ── Page de titre ──
\begin{titlepage}
\begin{center}
\vspace*{2cm}

{\Huge\bfseries\color{ifiBlue} O.S.C.A.R.}

\vspace{0.5cm}
{\Large\color{ifiGray} Outil de Suivi des Cours et d'Analyse du R\'{e}seau}

\vspace{1.5cm}
{\large Documentation technique et guide utilisateur}

\vspace{2cm}

\begin{tikzpicture}
  \draw[ifiBlue, line width=2pt] (0,0) -- (12,0);
\end{tikzpicture}

\vspace{1cm}
{\large\bfseries Institut fran\c{c}ais Italia}

\vspace{0.5cm}
{\color{ifiGray} R\'{e}seau des Instituts fran\c{c}ais en Italie\\
IFM (Milan) --- IFF (Florence) --- IFN (Naples) --- IFP (Palerme)}

\vspace{2cm}
{\color{ifiGray}\today}

\vspace{1cm}
{\small Version 2.8 --- Tableau de bord Streamlit}
\end{center}
\end{titlepage}

% ── Table des matières ──
\tableofcontents
\newpage

% =============================================================================
\section{Pr\'{e}sentation g\'{e}n\'{e}rale}
% =============================================================================

\subsection{Qu'est-ce qu'O.S.C.A.R.~?}

O.S.C.A.R. (\textbf{O}util de \textbf{S}uivi des \textbf{C}ours et d'\textbf{A}nalyse du \textbf{R}\'{e}seau) est un tableau de bord interactif d\'{e}velopp\'{e} pour l'Institut fran\c{c}ais Italia. Il permet d'analyser en d\'{e}tail l'activit\'{e} p\'{e}dagogique du r\'{e}seau \`{a} partir des donn\'{e}es export\'{e}es du logiciel de gestion AEC.

\subsection{Fonctionnalit\'{e}s principales}

\begin{itemize}[leftmargin=2em]
  \item \textbf{Analyse des cours par cat\'{e}gorie} --- r\'{e}partition par secteur, sous-secteur, macro-cat\'{e}gorie et cat\'{e}gorie
  \item \textbf{Fiches de cours} --- niveaux CECRL, inscriptions, tendances, nouveaux/r\'{e}inscrits
  \item \textbf{Profils clients} --- d\'{e}mographie, nationalit\'{e}s, motivations d'inscription, carte choropleth
  \item \textbf{Catalogue produits} --- types de produits, tarifs, grille horaire, analyse par antenne
  \item \textbf{Comparaisons} --- inter-antennes, inter-semestres, inter-annuelles
  \item \textbf{Carte d'Italie} --- inscriptions g\'{e}olocalis\'{e}es
  \item \textbf{Rentabilit\'{e} (ARPI)} --- revenu moyen par inscription
  \item \textbf{Export} --- Excel multi-feuilles et CSV
  \item \textbf{Aide rapide} --- assistant int\'{e}gr\'{e} pour l'exploration des donn\'{e}es
\end{itemize}

\subsection{Architecture technique}

\begin{center}
\begin{tabularx}{\textwidth}{lX}
\toprule
\textbf{Composant} & \textbf{D\'{e}tails} \\
\midrule
Langage & Python 3.9+ \\
Framework & Streamlit 1.30+ \\
Visualisation & Plotly Express, Plotly Graph Objects \\
Traitement & Pandas, NumPy \\
H\'{e}bergement & Streamlit Community Cloud \\
D\'{e}p\^{o}t & GitHub (\texttt{adrbn/ifi\_stats\_aec}) \\
Fichier principal & \texttt{dashboard\_aec\_v2.py} ($\sim$6\,900 lignes) \\
\bottomrule
\end{tabularx}
\end{center}

% =============================================================================
\section{Guide de d\'{e}marrage}
% =============================================================================

\subsection{Acc\`{e}s \`{a} l'application}

L'application est accessible en ligne \`{a} l'adresse :

\begin{center}
\url{https://ifi-stats-aec.streamlit.app}
\end{center}

\begin{infobox}[Connexion]
Deux comptes sont pr\'{e}configur\'{e}s. Les identifiants sont communiqu\'{e}s s\'{e}par\'{e}ment par le service informatique de l'IFI. La connexion est persistante gr\^{a}ce \`{a} un token HMAC stock\'{e} dans les param\`{e}tres d'URL.
\end{infobox}

\subsection{Pr\'{e}paration des donn\'{e}es}

Les donn\'{e}es proviennent exclusivement des exports du logiciel \textbf{AEC} (logiciel de gestion des cours des Alliances et Instituts fran\c{c}ais). Quatre types d'exports sont support\'{e}s :

\begin{enumerate}[leftmargin=2em]
  \item \textbf{Rapports par cat\'{e}gories} (\texttt{.xlsx}) --- un fichier par antenne et par semestre
  \item \textbf{Fiches de cours} (\texttt{.csv}) --- export global contenant les d\'{e}tails de chaque cours
  \item \textbf{Profils clients} (\texttt{.csv/.xlsx}) --- donn\'{e}es d\'{e}mographiques des inscrits
  \item \textbf{Catalogue produits} (\texttt{.xlsx}) --- r\'{e}f\'{e}rentiel des formations propos\'{e}es
\end{enumerate}

\begin{tipbox}[Conseil]
Utilisez le script d'export automatis\'{e} \texttt{aec\_exporter.py} (ou sa version GUI \texttt{aec\_exporter\_gui.py}) pour g\'{e}n\'{e}rer ces fichiers directement depuis AEC.
\end{tipbox}

\subsection{Chargement des fichiers}

\begin{enumerate}[leftmargin=2em]
  \item Cliquez sur \og{}Parcourir\fg{} dans la barre lat\'{e}rale gauche
  \item S\'{e}lectionnez un ou plusieurs fichiers (tous types confondus)
  \item O.S.C.A.R. d\'{e}tecte automatiquement le type de chaque fichier
  \item Les donn\'{e}es sont trait\'{e}es et les onglets d'analyse apparaissent
\end{enumerate}

\begin{infobox}[D\'{e}tection automatique]
Le syst\`{e}me identifie le type d'export (cat\'{e}gories, fiches, profils, produits) en analysant la structure des colonnes. Aucune configuration manuelle n'est n\'{e}cessaire.
\end{infobox}

% =============================================================================
\section{Onglets d'analyse}
% =============================================================================

\subsection{Synth\`{e}se (onglet principal)}

L'onglet de synth\`{e}se pr\'{e}sente les indicateurs-cl\'{e}s :

\begin{itemize}[leftmargin=2em]
  \item \textbf{Inscriptions} --- nombre total d'inscriptions sur la p\'{e}riode
  \item \textbf{Cours} --- nombre de cours programm\'{e}s (inclut \og{}Cours programm\'{e}s\fg{})
  \item \textbf{Heures-\'{e}l\`{e}ves} --- total des heures-\'{e}l\`{e}ves (inscriptions $\times$ heures)
  \item \textbf{Heures pr\'{e}vues} --- total des heures pr\'{e}vues de cours
  \item \textbf{Recettes} --- chiffre d'affaires total
  \item \textbf{\'{E}l\`{e}ves/cours} --- ratio moyen d'\'{e}l\`{e}ves par cours
\end{itemize}

\subsection{Par antenne}

Analyse ventil\'{e}e par antenne (IFM, IFF, IFN, IFP) avec :
\begin{itemize}[leftmargin=2em]
  \item Tableau r\'{e}capitulatif multi-indicateurs
  \item Graphiques en barres et camemberts
  \item Comparaisons entre antennes
\end{itemize}

\subsection{Par secteurs / sous-secteurs / macro-cat\'{e}gories / cat\'{e}gories}

Quatre niveaux de granularit\'{e} hierarchique :

\begin{center}
\begin{tabular}{lcl}
Secteur & $\rightarrow$ & Sous-secteur \\
 & $\rightarrow$ & Macro-cat\'{e}gorie \\
 & & \quad $\rightarrow$ Cat\'{e}gorie \\
\end{tabular}
\end{center}

Chaque niveau affiche tableau, graphiques et filtres d\'{e}di\'{e}s. Le fichier \texttt{data/category\_mapping.csv} d\'{e}finit les correspondances entre les cat\'{e}gories AEC brutes et la hi\'{e}rarchie secteur/sous-secteur/macro-cat\'{e}gorie.

\subsection{Comparaisons}

Trois modes de comparaison :
\begin{enumerate}[leftmargin=2em]
  \item \textbf{Antenne vs antenne} --- comparer deux sedi sur les m\^{e}mes indicateurs
  \item \textbf{Semestre vs semestre} --- S1 vs S2 (n\'{e}cessite les deux fichiers)
  \item \textbf{Secteur vs secteur} --- comparer deux secteurs d'activit\'{e}
\end{enumerate}

\subsection{Graphiques}

Navigation par carrousel dans une s\'{e}rie de visualisations :
\begin{itemize}[leftmargin=2em]
  \item Flux Sankey (antenne $\rightarrow$ secteur)
  \item Barres group\'{e}es (inscriptions par secteur et antenne)
  \item Treemap (r\'{e}partition hi\'{e}rarchique)
  \item Sunburst (vue en rayons de soleil)
\end{itemize}

\subsection{Fiches de cours}

Analyse des fiches de cours individuelles :
\begin{itemize}[leftmargin=2em]
  \item R\'{e}partition par niveau CECRL (A1, A2, B1, B2, C1, C2)
  \item Taux de remplissage des cours
  \item Nouveaux inscrits vs r\'{e}inscrits
  \item Synth\`{e}se par semestre et ann\'{e}e
\end{itemize}

\subsection{Profils clients}

Analyse d\'{e}mographique des inscrits :
\begin{itemize}[leftmargin=2em]
  \item \textbf{D\'{e}mographie} --- r\'{e}partition par \^{a}ge, sexe (M/F/Non sp\'{e}cifi\'{e})
  \item \textbf{Nationalit\'{e}s} --- carte choropleth mondiale, top nationalit\'{e}s
  \item \textbf{Motivations} --- raisons d'inscription, distribution
  \item \textbf{Cat\'{e}gories socio-professionnelles} --- CSP des inscrits
\end{itemize}

\subsection{Produits}

Analyse du catalogue de formations :
\begin{itemize}[leftmargin=2em]
  \item Types de produits et r\'{e}partition
  \item Grille tarifaire par type et famille
  \item Prix moyen par heure
  \item Comparaison entre antennes
\end{itemize}

\subsection{Export}

Deux formats disponibles :
\begin{itemize}[leftmargin=2em]
  \item \textbf{Excel multi-feuilles} (\texttt{.xlsx}) --- une feuille par antenne + totaux
  \item \textbf{CSV} (\texttt{.csv}) --- donn\'{e}es brutes combin\'{e}es
\end{itemize}

Des boutons \og{}T\'{e}l\'{e}charger CSV\fg{} sont \'{e}galement disponibles \`{a} la fin de chaque section (fiches, profils, produits).

\subsection{Aide rapide}

Module d'assistance int\'{e}gr\'{e} proposant :
\begin{itemize}[leftmargin=2em]
  \item Analyses automatiques (insights) sur les donn\'{e}es charg\'{e}es
  \item Questions pr\'{e}d\'{e}finies (total inscriptions, recettes, meilleure antenne\ldots)
  \item Champ libre pour poser des questions en langage naturel
\end{itemize}

% =============================================================================
\section{Donn\'{e}es et traitement}
% =============================================================================

\subsection{D\'{e}tection des fichiers}

Lors du chargement, O.S.C.A.R. proc\`{e}de \`{a} :

\begin{enumerate}[leftmargin=2em]
  \item \textbf{D\'{e}tection du format} --- CSV ou Excel, s\'{e}parateur auto-d\'{e}tect\'{e}
  \item \textbf{D\'{e}tection du type d'export} --- analyse des noms de colonnes pour identifier fiches, profils, produits ou cat\'{e}gories
  \item \textbf{Extraction des m\'{e}tadonn\'{e}es} --- antenne (sede), semestre et ann\'{e}e depuis le nom de fichier
\end{enumerate}

\subsection{Correspondance des cat\'{e}gories}

Le fichier \texttt{data/category\_mapping.csv} contient la hi\'{e}rarchie compl\`{e}te :

\begin{center}
\begin{tabularx}{\textwidth}{lX}
\toprule
\textbf{Colonne} & \textbf{Description} \\
\midrule
\texttt{Catégorie} & Nom de la cat\'{e}gorie AEC (exactement comme dans l'export) \\
\texttt{Secteur} & Secteur de rattachement (Langue, Culture, etc.) \\
\texttt{Sous-secteur} & Sous-secteur plus fin \\
\texttt{Macro-catégorie} & Regroupement synth\'{e}tique \\
\bottomrule
\end{tabularx}
\end{center}

\begin{warnbox}[Attention]
Toute cat\'{e}gorie absente de ce fichier appara\^{i}tra dans la section \og{}Cat\'{e}gories non rattach\'{e}es\fg{} de l'onglet Configuration. Il convient de mettre \`{a} jour le CSV r\'{e}guli\`{e}rement.
\end{warnbox}

\subsection{Cache de session}

Pour \'{e}viter de retraiter les m\^{e}mes fichiers \`{a} chaque interaction, O.S.C.A.R. utilise un cache bas\'{e} sur un hash MD5 des noms de fichiers :

\begin{lstlisting}[language=Python, caption={M\'{e}canisme de cache}]
_file_names_hash = hashlib.md5(
    "|".join(sorted(f.name for f in files)).encode()
).hexdigest()[:12]

if st.session_state.get('_files_hash') == _file_names_hash:
    # Restaurer depuis le cache
else:
    # Traitement complet + stockage en session
\end{lstlisting}

\subsection{Nationalit\'{e}s et g\'{e}olocalisation}

La correspondance nationalit\'{e} $\rightarrow$ code ISO-3 (pour la carte choropleth) utilise trois niveaux de correspondance :

\begin{enumerate}[leftmargin=2em]
  \item \textbf{Correspondance exacte} --- la nationalit\'{e} brute correspond exactement \`{a} une entr\'{e}e du dictionnaire
  \item \textbf{Correspondance par pr\'{e}fixe} --- les 6 premiers caract\`{e}res correspondent (seuil minimum pour \'{e}viter les faux positifs)
  \item \textbf{Correspondance floue} --- via \texttt{difflib.get\_close\_matches} avec un seuil de similarit\'{e} de 0.8
\end{enumerate}

% =============================================================================
\section{S\'{e}curit\'{e}}
% =============================================================================

\subsection{Authentification}

\begin{itemize}[leftmargin=2em]
  \item Les mots de passe sont stock\'{e}s sous forme de \textbf{hash SHA-256 statique} (aucun mot de passe en clair dans le code)
  \item En production, les identifiants peuvent \^{e}tre configur\'{e}s via \texttt{st.secrets} (fichier \texttt{.streamlit/secrets.toml})
  \item Le fallback local utilise des hashes pr\'{e}-calcul\'{e}s pour le d\'{e}veloppement
\end{itemize}

\subsection{Connexion persistante}

Un token HMAC (bas\'{e} sur le nom d'utilisateur et un secret) est stock\'{e} dans les param\`{e}tres d'URL. Cela permet de maintenir la session apr\`{e}s un rafra\^{i}chissement de page.

\begin{lstlisting}[language=Python, caption={G\'{e}n\'{e}ration du token HMAC}]
def _make_login_token(username):
    return hmac.new(
        _PERSIST_SECRET, username.encode(), hashlib.sha256
    ).hexdigest()[:20]
\end{lstlisting}

\begin{warnbox}[Important]
En production, le secret HMAC doit \^{e}tre d\'{e}fini dans \texttt{st.secrets} (cl\'{e} \texttt{hmac\_secret}) et ne pas utiliser la valeur par d\'{e}faut du code source.
\end{warnbox}

% =============================================================================
\section{D\'{e}ploiement}
% =============================================================================

\subsection{Streamlit Community Cloud}

L'application est d\'{e}ploy\'{e}e sur Streamlit Community Cloud. Chaque push sur la branche \texttt{main} d\'{e}clenche un red\'{e}ploiement automatique.

\subsubsection{Structure du d\'{e}p\^{o}t}

\begin{lstlisting}[language=bash, caption={Arborescence du d\'{e}p\^{o}t}]
ifi_stats_aec/
  dashboard_aec_v2.py    # Application principale
  requirements.txt       # Dependances Python
  data/
    category_mapping.csv  # Correspondances categories
    exports_AEC_*.zip     # Donnees pre-chargees
    export_*_produits_*.xlsx
    export_*_clients_*.csv
  img/
    logo/                 # Logos IFI
  docs/
    oscar_documentation.tex  # Ce document
\end{lstlisting}

\subsubsection{D\'{e}pendances}

\begin{lstlisting}[caption={requirements.txt}]
streamlit>=1.30
pandas>=2.0
plotly>=5.0
openpyxl
xlsxwriter
\end{lstlisting}

\subsubsection{Configuration des secrets}

Sur Streamlit Cloud, aller dans \textbf{Settings > Secrets} et ajouter :

\begin{lstlisting}[caption={.streamlit/secrets.toml}]
hmac_secret = "votre_secret_aleatoire_fort"

[users.adrien]
name = "Adrien"
password = "mot_de_passe"

[users.stephanie]
name = "Stephanie Sauvignon"
password = "mot_de_passe"
\end{lstlisting}

\subsection{Ex\'{e}cution locale}

\begin{lstlisting}[language=bash, caption={Lancement local}]
# Creer un environnement virtuel
python3 -m venv .venv
source .venv/bin/activate

# Installer les dependances
pip install -r requirements.txt

# Lancer l'application
streamlit run dashboard_aec_v2.py
\end{lstlisting}

% =============================================================================
\section{Param\'{e}trage avanc\'{e}}
% =============================================================================

\subsection{Codes des antennes}

\begin{center}
\begin{tabular}{lll}
\toprule
\textbf{Code} & \textbf{Ville} & \textbf{Institut} \\
\midrule
IFM & Milan & Institut fran\c{c}ais Milano \\
IFF & Florence & Institut fran\c{c}ais Firenze \\
IFN & Naples & Institut fran\c{c}ais Napoli \\
IFP & Palerme & Institut fran\c{c}ais Palermo \\
\bottomrule
\end{tabular}
\end{center}

\subsection{Conventions de nommage des fichiers}

Pour la d\'{e}tection automatique, les noms de fichiers doivent contenir :

\begin{itemize}[leftmargin=2em]
  \item Le code de l'antenne (\texttt{IFM}, \texttt{IFF}, \texttt{IFN}, \texttt{IFP})
  \item Optionnellement le semestre (\texttt{S1} ou \texttt{S2})
  \item Optionnellement l'ann\'{e}e (\texttt{2024}, \texttt{2025}\ldots)
\end{itemize}

Exemples valides :
\begin{itemize}[leftmargin=2em]
  \item \texttt{export\_IFM\_S1\_2024.xlsx}
  \item \texttt{rapport\_IFF\_annuel\_2025.xlsx}
  \item \texttt{fiches\_cours\_global.csv}
\end{itemize}

\subsection{Personnalisation CSS}

L'apparence du tableau de bord est enti\`{e}rement personnalis\'{e}e via la fonction \texttt{get\_css()} qui injecte du CSS dans la page Streamlit. Les principales personnalisations :

\begin{itemize}[leftmargin=2em]
  \item Masquage des \'{e}l\'{e}ments Streamlit natifs (toolbar, footer, hamburger menu)
  \item Bouton de collapse/expansion de la sidebar forc\'{e} en \texttt{position: fixed}
  \item Bouton d'impression PDF en position fixe (bas droite)
  \item Styles d'impression d\'{e}di\'{e}s (\texttt{@media print})
\end{itemize}

\subsection{Traductions}

O.S.C.A.R. est actuellement en fran\c{c}ais uniquement. Toutes les cha\^{i}nes de caract\`{e}res sont centralis\'{e}es dans le dictionnaire \texttt{TRANSLATIONS} et accessibles via la fonction \texttt{t(key)}. L'ajout d'une langue suppl\'{e}mentaire n\'{e}cessiterait d'ajouter un second dictionnaire et un s\'{e}lecteur de langue.

% =============================================================================
\section{Maintenance et mises \`{a} jour}
% =============================================================================

\subsection{Ajout d'une cat\'{e}gorie}

\begin{enumerate}[leftmargin=2em]
  \item Ouvrir \texttt{data/category\_mapping.csv}
  \item Ajouter une ligne avec : Cat\'{e}gorie, Secteur, Sous-secteur, Macro-cat\'{e}gorie
  \item Commiter et pusher sur \texttt{main}
  \item L'application se red\'{e}ploie automatiquement
\end{enumerate}

\subsection{Ajout d'une nationalit\'{e}}

Dans le dictionnaire \texttt{\_NATIONALITY\_TO\_ISO\_RAW} (ligne $\sim$2190 du fichier Python), ajouter une entr\'{e}e :

\begin{lstlisting}[language=Python]
_NATIONALITY_TO_ISO_RAW = {
    # ...
    "nouvelle nationalite": "ISO",  # Code ISO-3 du pays
}
\end{lstlisting}

\subsection{Modification des identifiants}

\begin{enumerate}[leftmargin=2em]
  \item Calculer le hash SHA-256 du nouveau mot de passe :
\begin{lstlisting}[language=Python]
import hashlib
hashlib.sha256("nouveau_mdp".encode()).hexdigest()
\end{lstlisting}
  \item Remplacer le hash correspondant dans la section fallback de \texttt{\_load\_users()}
  \item Ou mieux : configurer les identifiants dans \texttt{st.secrets} sur Streamlit Cloud
\end{enumerate}

% =============================================================================
\section{R\'{e}solution de probl\`{e}mes}
% =============================================================================

\begin{longtable}{p{5cm}p{9cm}}
\toprule
\textbf{Probl\`{e}me} & \textbf{Solution} \\
\midrule
\endhead
La sidebar ne se rouvre pas apr\`{e}s fermeture & V\'{e}rifier que le CSS pour \texttt{stSidebarCollapsedControl} est bien pr\'{e}sent (position fixed, z-index 999999). Rafra\^{i}chir la page. \\
\addlinespace
\og{}Aucune donn\'{e}e valide charg\'{e}e\fg{} & V\'{e}rifier le format des fichiers (CSV s\'{e}par\'{e} par point-virgule, Excel .xlsx). V\'{e}rifier que les colonnes attendues sont pr\'{e}sentes. \\
\addlinespace
Sede non d\'{e}tect\'{e}e & V\'{e}rifier que le nom du fichier contient bien IFM/IFF/IFN/IFP. \\
\addlinespace
Cat\'{e}gories non rattach\'{e}es & Mettre \`{a} jour \texttt{data/category\_mapping.csv} avec les nouvelles cat\'{e}gories. \\
\addlinespace
Nationalit\'{e}s non reconnues sur la carte & Ajouter la correspondance dans \texttt{\_NATIONALITY\_TO\_ISO\_RAW}. \\
\addlinespace
Erreur de d\'{e}ploiement Streamlit Cloud & V\'{e}rifier \texttt{requirements.txt} et la syntaxe Python du fichier principal. \\
\addlinespace
Connexion perdue apr\`{e}s rafra\^{i}chissement & V\'{e}rifier que les param\`{e}tres d'URL contiennent bien \texttt{user} et \texttt{token}. Reconfigurer le secret HMAC si n\'{e}cessaire. \\
\bottomrule
\end{longtable}

% =============================================================================
\section{Annexes}
% =============================================================================

\subsection{Liste des indicateurs}

\begin{center}
\begin{tabularx}{\textwidth}{lXl}
\toprule
\textbf{Indicateur} & \textbf{Description} & \textbf{Colonne source} \\
\midrule
Inscriptions & Nombre total d'inscriptions & \texttt{Nb. d'inscriptions} \\
Cours & Nombre de cours (y.c. programm\'{e}s) & \texttt{Cours programm\'{e}s} \\
Heures-\'{e}l\`{e}ves & Inscriptions $\times$ heures & \texttt{Total des heures-\'{e}l\`{e}ves} \\
Heures pr\'{e}vues & Heures planifi\'{e}es & \texttt{Total des heures pr\'{e}vues} \\
Recettes & Chiffre d'affaires & \texttt{Total g\'{e}n\'{e}ral TTC} \\
\'{E}l\`{e}ves/cours & Ratio moyen & Inscriptions / Cours \\
ARPI & Average Revenue Per Inscription & Recettes / Inscriptions \\
\bottomrule
\end{tabularx}
\end{center}

\subsection{Technologies utilis\'{e}es}

\begin{center}
\begin{tabular}{ll}
\toprule
\textbf{Biblioth\`{e}que} & \textbf{Usage} \\
\midrule
Streamlit & Framework web interactif \\
Pandas & Manipulation de donn\'{e}es \\
Plotly & Visualisations interactives \\
OpenPyXL & Lecture de fichiers Excel \\
XlsxWriter & \'{E}criture de fichiers Excel \\
difflib & Correspondance floue (nationalit\'{e}s) \\
hashlib & Hachage SHA-256, MD5 \\
hmac & Tokens d'authentification persistants \\
\bottomrule
\end{tabular}
\end{center}

\vfill
\begin{center}
\color{ifiGray}\small
Document g\'{e}n\'{e}r\'{e} le \today{} --- O.S.C.A.R. v2.8\\
Institut fran\c{c}ais Italia
\end{center}

\end{document}
